% =============================================================================
%  CHAPTER 4 -- REGULARIZATION: RIDGE, LASSO, ELASTIC NET, PCR
% =============================================================================
\section{Regularization: ridge, lasso, elastic net, principal components}
\label{ch:regularization}

Chapter~\ref{ch:ols} ended with OLS being unbiased but high-variance when regressors are
correlated or numerous. This chapter derives the four remedies the project implements,
in C++ (\code{src/model/basket\_selector.h}) and in Python (sklearn), and explains what
each one's geometry implies for a hedging problem.

\subsection{The general idea: constrained estimation}
\label{sec:reg:idea}

Every regularization method solves
\begin{equation}
  \hat\beta \;=\; \argmin_{\beta}\ \Bigl\{\ \underbrace{\tfrac12\lVert Y-X\beta\rVert_2^2}_{\text{fit}}
  \;+\; \underbrace{\mathrm{Pen}(\beta)}_{\text{complexity}}\ \Bigr\},
  \label{eq:reg:general}
\end{equation}
with $\mathrm{Pen}$ chosen to encode a prior belief about what a reasonable $\beta$ looks
like. By the Lagrangian duality of convex problems, the ``penalty'' form
\eqref{eq:reg:general} is equivalent to a ``constraint'' form
$\min_\beta \tfrac12\lVert Y-X\beta\rVert^2$ subject to $\mathrm{Pen}(\beta)\le s$, with
$\lambda$ and $s$ in one-to-one correspondence. The constraint picture is more
illuminating: \emph{you are choosing the best-fitting point inside a shrunken region of
parameter space}, and the shape of that region determines the qualitative behaviour.

\begin{intuition}
An L2 constraint region is a ball; an L1 region is a diamond (an octahedron in higher
dimensions). The fit's contour ellipses are typically very elongated along the
collinear directions. An elongated ellipse touching a \emph{ball} lands somewhere on the
surface with all coordinates non-zero; touching a \emph{diamond} overwhelmingly lands on
a \emph{corner}, and corners are where coordinates are exactly zero. That single picture
is the entire difference between ridge (shrinks) and lasso (selects).
\end{intuition}

\subsection{Ridge regression}
\label{sec:ridge}

\subsubsection{Derivation}
\label{sec:ridge:derivation}

Take $\mathrm{Pen}(\beta)=\tfrac{\lambda}{2}\lVert\beta\rVert_2^2$ with $\lambda\ge0$:
\begin{equation}
  \hat\beta_{\text{ridge}} = \argmin_\beta\ \Bigl\{\tfrac12 (Y-X\beta)^{\top}(Y-X\beta)
  + \tfrac{\lambda}{2}\beta^{\top}\beta\Bigr\}.
\end{equation}
The gradient is $-X^{\top}(Y-X\beta)+\lambda\beta$; setting it to zero,
\begin{equation}
  (X^{\top}X+\lambda I)\,\hat\beta_{\text{ridge}} \;=\; X^{\top}Y
  \quad\Longrightarrow\quad
  \boxed{\ \hat\beta_{\text{ridge}} = \bigl(X^{\top}X+\lambda I\bigr)^{-1}X^{\top}Y\ }
  \label{eq:ridge}
\end{equation}
Since $X^{\top}X+\lambda I$ has eigenvalues $\lambda_i+\lambda>0$, the solution always
exists and is unique, even for a singular $X^{\top}X$ --- which is why ridge is also a
numerical stabiliser (Section~\ref{sec:ols:multicol}).

\subsubsection{Bias and variance, exactly}
\label{sec:ridge:biasvar}

Write $A_\lambda=(X^{\top}X+\lambda I)^{-1}X^{\top}$ and
$M = (X^{\top}X+\lambda I)^{-1}X^{\top}X$. Then
$\hat\beta_{\text{ridge}} = M\beta + A_\lambda\eps$, so
\begin{align}
  \operatorname{Bias}(\hat\beta_{\text{ridge}}) &= (M-I)\beta
    = -\lambda\bigl(X^{\top}X+\lambda I\bigr)^{-1}\beta, \label{eq:ridge:bias}\\
  \Var(\hat\beta_{\text{ridge}}) &= \sigma^2 A_\lambda A_\lambda^{\top}
    = \sigma^2\bigl(X^{\top}X+\lambda I\bigr)^{-1}X^{\top}X\bigl(X^{\top}X+\lambda I\bigr)^{-1}.
    \label{eq:ridge:var}
\end{align}
Derivation of \eqref{eq:ridge:bias}: $M-I=(X^{\top}X+\lambda I)^{-1}(X^{\top}X - X^{\top}X-\lambda I)
= -\lambda(X^{\top}X+\lambda I)^{-1}$. As $\lambda\to0$ the bias vanishes and
\eqref{eq:ridge:var} becomes the OLS variance; as $\lambda\to\infty$ the estimator
collapses to $0$ (zero variance, maximal bias). The MSE is minimized at an interior
$\lambda^\ast>0$ whenever $\lVert\beta\rVert$ is finite and $n\ge2$ --- the celebrated
fact that \emph{a biased estimator can beat OLS on MSE}. (Hoerl--Kennard existence
result; the unit test in \code{test/engine\_tests.cpp} asserts the
$\lambda\to0$ limit numerically.)

\subsubsection{The spectral view: ridge is a low-pass filter}
\label{sec:ridge:spectral}

Use $X^{\top}X=V\Lambda V^{\top}$ (Section~\ref{sec:prob:linalg}) and rotate coordinates:
$\tilde\beta=V^{\top}\beta$, $\tilde Z=V^{\top}X^{\top}Y$. Then
\begin{equation}
  \hat{\tilde\beta}_{\text{ridge},i} \;=\; \frac{\lambda_i}{\lambda_i+\lambda}\,
  \hat{\tilde\beta}_{\text{OLS},i},
  \label{eq:ridge:spectral}
\end{equation}
i.e.\ each principal direction of the regressor space is multiplied by a shrinkage
factor $d_i=\lambda_i/(\lambda_i+\lambda)\in(0,1]$. Directions with large eigenvalues
(strongly identified by the data) pass through almost unchanged; directions with small
eigenvalues (weakly identified, where OLS is most volatile) are damped toward zero. This
is exactly a \emph{low-pass filter} in eigenspace, and it explains both the empirical
findings of Phase~\ref{ch:phase5} (ridge reduces hedge-weight churn: coefficient
instability $0.998\to0.935$) and the numerical benefit (condition number bounded by
$(\lambda_{\max}+\lambda)/\lambda$).

\subsubsection{The Bayesian view: ridge is a Gaussian prior}
\label{sec:ridge:bayes}

Let $Y\mid\beta\sim\mathcal N(X\beta,\sigma^2 I)$ and put a prior
$\beta\sim\mathcal N(0,\tau^2 I)$. The posterior is
\[
  p(\beta\mid Y) \;\propto\;
  \exp\!\Bigl(-\frac{1}{2\sigma^2}\lVert Y-X\beta\rVert^2\Bigr)
  \exp\!\Bigl(-\frac{1}{2\tau^2}\lVert\beta\rVert^2\Bigr),
\]
so the log-posterior is $-\frac{1}{2\sigma^2}\{\lVert Y-X\beta\rVert^2 +
\frac{\sigma^2}{\tau^2}\lVert\beta\rVert^2\}$ and its mode is
\eqref{eq:ridge} with $\lambda=\sigma^2/\tau^2$. Because both prior and likelihood are
Gaussian, the posterior is Gaussian and the mode equals the mean: \emph{ridge is the
posterior mean under a zero-mean Gaussian prior}. The prior variance $\tau^2$ is the
knob: believing a priori that hedge weights are small ($\tau$ small) is the same as
penalising them hard ($\lambda$ large). Economically: ``unless the data insists
otherwise, hold a small, well-conditioned hedge.''

\subsubsection{Generalised cross-validation and leave-one-out}
\label{sec:ridge:loo}

How do you choose $\lambda$ honestly? Cross-validation: hold out data, fit on the rest,
score the held-out prediction error, pick the $\lambda$ that minimises it. For linear
smoothers ($\hat Y = S_\lambda Y$ with $S_\lambda = X(X^{\top}X+\lambda I)^{-1}X^{\top}$)
the leave-one-out error has a closed form that needs \emph{one} fit, not $T$:
\begin{equation}
  \mathrm{LOO}(\lambda) \;=\; \frac{1}{T}\sum_{t=1}^{T}
  \Bigl(\frac{y_t-\hat y_t}{1-[S_\lambda]_{tt}}\Bigr)^{2},
  \qquad
  \mathrm{GCV}(\lambda) \;=\; \frac{\mathrm{RSS}(\lambda)/T}{\bigl(1-\tr(S_\lambda)/T\bigr)^{2}} .
  \label{eq:gcv}
\end{equation}
The derivation of the LOO identity uses Sherman--Morrison \eqref{eq:shermanmorrison} on
the rank-1 removal of observation $t$ from $X^{\top}X$. GCV replaces the individual
leverages by their average and is the standard cheap surrogate.

\begin{remark}[Honesty: we do not use GCV here]
Phase~\ref{ch:phase5} selects $\lambda$ by \emph{walk-forward net Sharpe on a validation
window} (2019--2021), not by prediction error, because the quantity we care about is the
traded P\&L after costs, not the fit. That choice has a cost: it makes the selection
itself a hypothesis that must be counted in the multiple-testing audit
(Phase~\ref{ch:phase10} counts it: $M=44$). Selecting on prediction error would be more
conservative statistically and less relevant economically. Both choices are defensible;
the important thing is that the selection window is disjoint from the reported test
window.
\end{remark}

\subsection{Lasso}
\label{sec:lasso}

\subsubsection{The problem and why there is no closed form}

Take $\mathrm{Pen}(\beta)=\lambda\lVert\beta\rVert_1=\lambda\sum_j\lvert\beta_j\rvert$:
\begin{equation}
  \hat\beta_{\text{lasso}} = \argmin_\beta\ \Bigl\{\tfrac12\lVert Y-X\beta\rVert_2^2
  + \lambda\lVert\beta\rVert_1\Bigr\}.
  \label{eq:lasso}
\end{equation}
The absolute value is not differentiable at zero, so we work with \emph{subdifferentials}.
For a convex $f$, $\partial f(\beta)$ is the set of subgradients $g$ with
$f(b)\ge f(\beta)+g^{\top}(b-\beta)$ for all $b$. For $\lvert\beta_j\rvert$,
\[
  \partial\lvert\beta_j\rvert =
  \begin{cases}
    \{+1\} & \beta_j>0,\\
    \{-1\} & \beta_j<0,\\
    [-1,+1] & \beta_j=0 .
  \end{cases}
\]
A point $\hat\beta$ minimises \eqref{eq:lasso} iff $0\in\partial$ of the objective, i.e.\
for every $j$
\begin{equation}
  x_j^{\top}(Y-X\hat\beta) \;=\; \lambda\, s_j,
  \qquad
  s_j \in \partial\lvert\hat\beta_j\rvert .
  \label{eq:lasso_kkt}
\end{equation}
Read \eqref{eq:lasso_kkt} carefully; it is the whole story of sparsity:
\begin{itemize}
  \item If $\hat\beta_j\ne0$, then $s_j=\sign(\hat\beta_j)$ and the residual correlation
    $x_j^{\top}\hat\eps$ equals $\pm\lambda$ exactly: every included variable is
    ``pushing'' with the maximum force the penalty allows.
  \item If $\hat\beta_j=0$, the condition is only
    $\lvert x_j^{\top}\hat\eps\rvert\le\lambda$: a variable stays out as long as its
    marginal contribution to the fit is \emph{smaller than the penalty threshold}. This
    is a genuine inequality, so solutions with $\hat\beta_j=0$ occur on a set of nonzero
    measure --- sparsity is not an accident, it is the generic outcome.
\end{itemize}
This is impossible with the L2 penalty, whose derivative at zero is zero, so the
equivalent condition would force $x_j^{\top}\hat\eps=0$ exactly --- a measure-zero event.

\subsubsection{Coordinate descent and the soft-threshold operator}
\label{sec:lasso:cd}

Because \eqref{eq:lasso_kkt} cannot be solved in closed form for general $X$, we minimise
cyclically over one coordinate at a time. Fixing all $\beta_k$, $k\ne j$, the objective as
a function of $\beta_j$ alone is
\[
  \tfrac12\lVert r_{-j} - X_j\beta_j\rVert^2 + \lambda\lvert\beta_j\rvert,
  \qquad r_{-j}=Y-\sum_{k\ne j}X_k\beta_k \ \text{(the partial residual)}.
\]
Expanding, this is a one-dimensional quadratic plus an absolute value; minimising over
$\beta_j>0$ and $\beta_j<0$ separately and comparing gives the closed-form update
\begin{equation}
  \hat\beta_j \;\leftarrow\; \frac{1}{\rho_j}\,S_{\lambda}\!\bigl(X_j^{\top}r_{-j}\bigr),
  \qquad \rho_j = X_j^{\top}X_j,
  \qquad
  S_{\lambda}(z)=\sign(z)\bigl(\lvert z\rvert-\lambda\bigr)_{+},
  \label{eq:softthresh}
\end{equation}
where $(u)_+=\max(u,0)$. $S_\lambda$ is the \emph{soft-threshold} (or shrinkage)
operator: subtract $\lambda$ from the magnitude, keep the sign, and clip at zero.
Iterate over $j=1,\dots,n$ until no coordinate changes by more than a tolerance; sweep
$\lambda$ over a decreasing path (warm-starting from the previous, sparser solution) to
get the whole lasso path in a few passes. Coordinate descent is guaranteed to converge
here because the objective is convex and \emph{separable} in the non-smooth part --- the
standard sufficient condition (Tseng's convergence result for block-coordinate descent on
smooth-plus-separable-convex objectives).

\paragraph{Standardisation matters.} The penalty $\lambda\lvert\beta_j\rvert$ depends on
the units of $X_j$: a regressor measured in dollars-per-share is penalised differently
from one measured in log prices. All lasso results in this book are computed on
column-standardised features ($x_j\mapsto(x_j-\bar x_j)/s_j$), with the intercept
excluded from the penalty, and the coefficients mapped back afterwards. Not doing this is
a common and silently damaging error.

\subsection{Elastic net}
\label{sec:elasticnet}

\begin{equation}
  \hat\beta_{\text{enet}} = \argmin_\beta\ \Bigl\{\tfrac12\lVert Y-X\beta\rVert^2
  + \lambda_2\lVert\beta\rVert_2^2 + \lambda_1\lVert\beta\rVert_1\Bigr\},
  \qquad
  \text{often written } \lambda\bigl[\rho\lVert\beta\rVert_1 + \tfrac{1-\rho}{2}\lVert\beta\rVert_2^2\bigr].
  \label{eq:elasticnet}
\end{equation}
Its motivation is a specific lasso pathology: with $p$ highly correlated regressors, lasso
tends to pick \emph{one} of them arbitrarily (whichever the data path hits first) and
ignore the rest, and the selection is unstable under small data perturbations. The L2
part couples correlated coefficients --- the penalty is minimised by \emph{equal} weights
on perfectly correlated regressors rather than by one large weight --- giving a
``grouping effect''. The coordinate-descent update is the same soft-threshold with an
inflated denominator:
\begin{equation}
  \hat\beta_j \leftarrow \frac{S_{\lambda_1}\bigl(X_j^{\top}r_{-j}\bigr)}{\rho_j+2\lambda_2},
  \label{eq:enet:cd}
\end{equation}
so elastic net is a one-line change to a lasso implementation. Equivalently, elastic net
is a lasso on the augmented data $\tilde X=[X;\sqrt{\lambda_2}I]$, $\tilde Y=[Y;0]$.

\subsection{Principal components and principal component regression}
\label{sec:pcr}

\subsubsection{Deriving PCA as a variance-maximisation problem}
\label{sec:pca:derivation}

Centre the regressors: $\tilde X = X-\mathbf 1\bar x^{\top}$, and let
$S=\tfrac1T \tilde X^{\top}\tilde X$ be the sample covariance. We seek the direction
$v$ (unit norm) along which the projected data $\tilde X v$ has the largest variance:
\begin{equation}
  \max_{v}\ v^{\top}Sv \quad\text{s.t.}\quad v^{\top}v=1 .
  \label{eq:pca:max}
\end{equation}
Lagrangian: $\mathcal L = v^{\top}Sv-\lambda(v^{\top}v-1)$, gradient
$2Sv-2\lambda v=0$, hence
\begin{equation}
  S v = \lambda v ,
  \label{eq:pca:eigen}
\end{equation}
so the candidate directions are eigenvectors, and the objective value at such a $v$ is
$v^{\top}Sv=\lambda$. Maximising means taking the \emph{largest} eigenvalue. The second
component is the maximiser subject to being orthogonal to the first, which yields the
second-largest eigenvalue, and so on. The spectral theorem
\eqref{eq:spectral} then gives $S=V\Lambda V^{\top}$ with
$\lambda_1\ge\lambda_2\ge\dots\ge\lambda_n\ge0$ and orthonormal columns $v_i$.

Components and scores: $Z=\tilde X V\in\R^{T\times n}$, with
$\Cov(Z)=\Lambda$ (diagonal: components are uncorrelated) and
$\Var(Z_{\cdot k})=\lambda_k$. The \emph{explained variance ratio} of component $k$ is
$\lambda_k/\sum_i\lambda_i$, and the cumulative ratio tells you how many components to
keep.

\begin{intuition}
In a basket of same-sector equities, $\lambda_1$ is huge relative to the rest: the first
principal component is the common sector/market factor, and $v_1$ is roughly proportional
to each name's volatility. Components $2,3,\dots$ are relative-value modes (``software
versus semis within tech'', ``Adobe versus Autodesk''). A stat-arb trader cares about the
\emph{trailing} components --- the relative-value modes --- while a risk manager cares
about the \emph{leading} ones. This tension explains the PCR result of
Phase~\ref{ch:phase5}: keeping only component 1 removes the market mode from the residual
and produces the \emph{best raw signal} (highest IC, $0.21$) but a book that under-hedges
the names' common drift and therefore the \emph{worst realized net Sharpe} ($0.03$).
\end{intuition}

\subsubsection{PCR and its relation to ridge}

Regress $Y$ on the first $k$ components:
\begin{equation}
  \hat\beta_{\text{PCR}} \;=\; V_{:k}\bigl(Z_{:k}^{\top}Z_{:k}\bigr)^{-1}Z_{:k}^{\top}Y
  \;=\; \sum_{i=1}^{k}\frac{v_i^{\top}\tilde X^{\top}Y}{T\lambda_i}\,v_i .
  \label{eq:pcr}
\end{equation}
(The second equality uses $Z_{:k}^{\top}Z_{:k}=T\Lambda_k$.) Compare with ridge in the
same rotated basis, equation~\eqref{eq:ridge:spectral}: ridge multiplies component $i$ by
$\lambda_i/(\lambda_i+\lambda)$ --- a \emph{smooth} shrinkage --- whereas PCR multiplies
by $1$ if $i\le k$ and $0$ otherwise --- a \emph{hard} cutoff. PCR is ridge with a step
function; ridge is PCR with a smooth one. Since the eigenvalue ordering is only loosely
related to how predictive each direction is, ridge is usually the safer default, which is
what Phase~\ref{ch:phase5} finds.

\subsubsection{Computing the eigendecomposition without a library}
\label{sec:pca:jacobi}

The sandbox has no Eigen and no LAPACK, so \code{src/model/dense\_la.h} implements a
\emph{cyclic Jacobi} eigen-solver: repeatedly pick a pair of coordinates $(p,q)$, compute
the rotation angle that annihilates the off-diagonal element $A_{pq}$,
\begin{equation}
  \tau = \frac{A_{qq}-A_{pp}}{2A_{pq}},\qquad
  t = \frac{\sign(\tau)}{\lvert\tau\rvert+\sqrt{1+\tau^2}},\qquad
  c = \frac{1}{\sqrt{1+t^2}},\quad s = tc,
  \label{eq:jacobi}
\end{equation}
apply the Givens rotation $A\leftarrow J^{\top}AJ$ (with $J$ the identity except
$J_{pp}=J_{qq}=c$, $J_{pq}=-s$, $J_{qp}=s$), and accumulate $V\leftarrow VJ$. Cycling over
all pairs reduces the off-diagonal Frobenius norm monotonically; convergence is quadratic
in the final sweeps. Cost $O(n^3)$ per sweep, which for $n\le5$ is nothing. Jacobi is
chosen over Householder tridiagonalisation + QR because it is (i) about 60 lines with no
scratch memory, (ii) extremely accurate for small symmetric matrices (it computes
eigenvalues to high relative accuracy even when they are clustered), and (iii) trivially
unit-testable --- the test asserts $V^{\top}V=I$, $AV=V\Lambda$, and that
$\mathrm{PCR}(k{=}n)$ reproduces OLS exactly.

\subsection{Choosing among the four: what the geometry implies for hedging}
\label{sec:reg:comparison}

\begin{center}
\renewcommand{\arraystretch}{1.25}
\begin{tabular}{p{0.13\linewidth}p{0.24\linewidth}p{0.24\linewidth}p{0.28\linewidth}}
\toprule
Method & Penalty / constraint & Effect on weights & When it is the right prior for a \emph{hedge} \\
\midrule
OLS & none & unbiased, high variance under collinearity & small basket, low collinearity, ample data \\
Ridge & $\lambda\lVert\beta\rVert_2^2$ & all shrunk smoothly, none zero, correlated names share weight & \textbf{hedging}: every name carries shared factor exposure, so you want them all, just not wildly \\
Lasso & $\lambda\lVert\beta\rVert_1$ & sparse; exactly zero beyond a threshold & \emph{selection}: large candidate pool where most names are irrelevant \\
Elastic net & $\lambda_1\lVert\beta\rVert_1+\lambda_2\lVert\beta\rVert_2^2$ & sparse but groups correlated names & selection among many near-duplicates \\
PCR & hard cutoff in eigenspace & weight only on top-$k$ directions & strong single factor, few effective dimensions \\
\bottomrule
\end{tabular}
\end{center}

The key insight, confirmed empirically in Phase~\ref{ch:phase5}: \emph{lasso's sparsity
is the wrong inductive bias for a hedging regression}. Sparsity is what you want when the
truth is that most candidate variables do not matter. In a three-name software basket, all
three names matter --- INTU's individual coefficient is statistically insignificant over
the full sample (HAC $t\approx0.45$, Table~\ref{tab:diag_coint}) yet removing it from the
hedge still degrades the book, because it reduces residual variance in the rolling
estimate. With $n=3$ and $\lambda$ chosen on validation, lasso and elastic net both
collapse the hedge, and the net Sharpe falls from $0.050$ (OLS) to $0.007$ (lasso) and
$-0.032$ (elastic net). Ridge, at $\alpha=10$, raises it to $0.094$.

\subsection{Model selection discipline}
\label{sec:reg:selection}

Three rules, all applied in Part~II:
\begin{enumerate}
  \item \textbf{Nested windows.} train 2010--2018 / validation 2019--2021 / test
    2022--2026, with the test period touched exactly once, at the end, for reporting.
    Hyper-parameters ($\lambda$, $k$, window $W$, $z_{\text{entry}}$) are chosen on
    validation only.
  \item \textbf{Small coarse grids.} Ridge $\alpha\in\{0.1,1,10\}$, lasso/enet
    $\alpha\in\{10^{-4},10^{-3},10^{-2}\}$, PCR $k\in\{1,2,3\}$, window
    $W\in\{60,120,180,252\}$. Coarse grids limit the search space that the
    multiple-testing audit must account for --- every grid point is a hypothesis.
  \item \textbf{Count everything.} The audit of Phase~\ref{ch:phase10} enumerates
    $M=44$ strategies actually examined, and the corrections are computed with that $M$.
    A grid of 12 hyper-parameter settings and 4 baskets is not ``free''.
\end{enumerate}
